\section*{Bab \ref{ch:g12:internal-heat} --- Panas Internal Bumi}

\begin{solution}{exo:g12:internal-heat:1}
(a) \omterm{def:g12:internal-heat:heat-flow}{aliran panas}; (b) \omterm{def:g12:internal-heat:heat-production}{produksi panas radiogenik}; (c) \omterm{def:g12:internal-heat:secular-cooling}{pendinginan sekuler};
(d) \omterm{def:g12:internal-heat:convection}{konveksi termal}; (e) \omterm{def:g12:internal-heat:geothermal}{energi panas bumi}.
\end{solution}

\begin{solution}{exo:g12:internal-heat:2}
(1) \qty{0.0916}{W/m^2}. (2) \qty{4.67e13}{W}. (3) \qty{0.029}{K/m}: selisih
satu derajat Celsius sama dengan selisih satu kelvin.
\end{solution}

\begin{solution}{exo:g12:internal-heat:3}
Torium-232, uranium-238, kalium-40, dan uranium-235. KamLAND hanya melihat
geoneutrino torium-232 dan uranium-238.
\end{solution}

\begin{solution}{exo:g12:internal-heat:4}
$490/\sqrt{1} = \qty{490}{mW/m^2}$; $490/2 = \qty{245}{mW/m^2}$;
$490/5 = \qty{98}{mW/m^2}$; $490/7 = \qty{70}{mW/m^2}$, kira-kira rata-rata
benua.
\end{solution}

\begin{solution}{exo:g12:internal-heat:5}
Di dekat permukaan, suhu tanah mengikuti hari dan musim; baru di bawah
beberapa meter suhu itu tetap pada suhu rata-rata tahunan, dan baru lebih dalam
lagi suhu itu menunjukkan panas yang datang dari bawah. Pengeboran mengaduk
batuan dan air di dalam lubang dan mengubah suhunya, sehingga kita menunggu
sampai lubang itu kembali ke suhu batuannya.
\end{solution}

\begin{solution}{exo:g12:internal-heat:6}
(1) $(265 - 10)/9.1 \approx \qty{28.0}{\celsius}$ per kilometer, yaitu
\qty{0.0280}{K/m}. (2) $q = 2.72 \times 0.0280 \approx \qty{0.0762}{W/m^2}$,
sekitar \qty{76}{mW/m^2}. (3) Sedikit di atas rata-rata benua,
\qty{70.9}{mW/m^2}, sekitar \qty{7}{\percent}: nilai benua yang umum.
\end{solution}

\begin{solution}{exo:g12:internal-heat:7}
Benua: $0.0709 \times \num{2.073e14} \approx \qty{1.47e13}{W} =
\qty{14.7}{TW}$. Samudra: $0.1054 \times \num{3.028e14} \approx
\qty{3.19e13}{W} = \qty{31.9}{TW}$. Jumlahnya sekitar \qty{46.6}{TW}, dan
$31.9/46.6 \approx \qty{68}{\percent}$ di antaranya, dua pertiga, keluar melalui
dasar samudra.
\end{solution}

\begin{solution}{exo:g12:internal-heat:8}
(1) $\num{1.88e-6} \times 8000 \approx \qty{0.0150}{W/m^2} = \qty{15.0}{mW/m^2}$.
(2) $15.0/70.9 \approx \qty{21}{\percent}$, seperlima \omterm{def:g12:internal-heat:heat-flow}{aliran panas}.
\end{solution}

\begin{solution}{exo:g12:internal-heat:9}
$\sqrt{t} = 490/70.9 \approx 6.91$, sehingga $t \approx 47.8$ juta tahun: dasar
samudra yang lebih tua dari sekitar 48 juta tahun kehilangan lebih sedikit panas
daripada benua rata-rata.
\end{solution}

\begin{solution}{exo:g12:internal-heat:10}
Uranium-235: $2^{1.4/0.704} = 2^{1.99} \approx 4$ kali lebih banyak. Kalium-40:
$2^{2.5/1.248} = 2^{2.00} \approx 4$ kali lebih banyak. Keduanya berjarak dua
waktu paruh ke belakang.
\end{solution}

\begin{solution}{exo:g12:internal-heat:11}
(1) $10^2/10^{-6} = \qty{1e8}{s}$, sekitar 3 tahun. (2) $(10^5)^2/10^{-6} =
\qty{1e16}{s}$, sekitar 320 juta tahun. (3) $(\num{2.891e6})^2/10^{-6}
\approx \qty{8.4e18}{s}$, sekitar 265 miliar tahun. Waktunya bertambah sebanding
dengan kuadrat jaraknya: sepuluh kali lebih tebal, seratus kali lebih lama.
\end{solution}

\begin{solution}{exo:g12:internal-heat:12}
Kalikan masing-masing dengan $2^{3/T}$: torium-232, $10.4 \times 2^{3/14.0} \approx
\qty{12.1}{TW}$; uranium-238, $9.6 \times 2^{3/4.468} \approx \qty{15.3}{TW}$;
kalium-40, $4.0 \times 2^{3/1.248} \approx \qty{21.2}{TW}$; uranium-235,
$0.42 \times 2^{3/0.704} \approx \qty{8.1}{TW}$. Jumlahnya sekitar
\qty{56.7}{TW}, lebih dari dua kali \qty{24.4}{TW} saat ini. Kedua isotop yang
berumur pendek memberikan $(21.2 + 8.1)/56.7 \approx \qty{52}{\percent}$ dari
jumlah itu, dibandingkan dengan $4.42/24.4 \approx \qty{18}{\percent}$ saat ini.
\end{solution}

\begin{solution}{exo:g12:internal-heat:13}
Daya adalah massa per detik dikalikan kalor jenis dikalikan kenaikan suhu:
$\num{3e8} = 1640 \times 4180 \times \Delta T$, sehingga $\Delta T =
\num{3e8}/(1640 \times 4180) \approx \qty{44}{K}$. Benua rata-rata kehilangan
\qty{0.0709}{W} per meter persegi, sehingga \qty{300}{MW} adalah kehilangan
panas alami seluas $\num{3e8}/0.0709 \approx \qty{4.2e9}{m^2}$, sekitar
\qty{4200}{km^2}: pembangkit itu mengumpulkan panas yang secara alami dilepaskan
Bumi di atas ribuan kilometer persegi.
\end{solution}

\begin{solution}{exo:g12:internal-heat:14}
(1) 100 juta tahun adalah $\num{2.6e11}/\num{1e8} = 2600$ kali lebih singkat
daripada 260 miliar tahun; 200 juta tahun 1300 kali lebih singkat.
(2) Konveksi mengangkut panas bersama batuan yang bergerak: waktu untuk
melintasi mantel adalah tebalnya dibagi kecepatannya, $\qty{2891}{km}$ dengan
laju beberapa sentimeter setahun, sekitar seratus juta tahun. Konduksi tidak
mempunyai pengangkut seperti itu: waktunya bertambah sebanding dengan kuadrat
jarak, dan di sepanjang ribuan kilometer waktunya menjadi sangat besar.
\end{solution}

\begin{solution}{exo:g12:internal-heat:15}
(1) $10/22.3 \approx \qty{45}{\percent}$; sisanya adalah pendinginan mantel
itu sendiri, yang kehilangan panas simpanannya sendiri. (2) Batas itu terletak pada jari-jari $6371 - 2891 =
\qty{3480}{km}$; luasnya $4\pi (\num{3.48e6})^2 \approx
\qty{1.52e14}{m^2}$, dan $\num{1e13}/\num{1.52e14} \approx
\qty{0.066}{W/m^2}$, sekitar \qty{66}{mW/m^2}, kira-kira \qty{72}{\percent}
dari rata-rata di permukaan. (3) Panas melintasi batas itu secara konduksi ke
dalam lapisan mantel yang paling bawah; setelah menghangat, batuan ini menjadi
lebih ringan daripada mantel di atasnya dan dapat naik, sambil mengangkut panas
ke atas: dengan cara inilah gumpalan diduga bermula, jauh di bawah.
\end{solution}

\begin{solution}{pb:g12:internal-heat:1}
\textbf{1.} $4\pi (\num{6.371e6})^2 \approx \qty{5.101e14}{m^2}$, dan
$2.073 + 3.028 = 5.101$: ya.

\textbf{2.} $3.028/5.101 \approx \qty{59}{\percent}$, sekitar tiga perlima.

\textbf{3.} Dasar samudra itu muda: dasar samudra dibuat dalam keadaan panas
di punggungan dan mendingin ketika bergerak menjauh, dan \omterm{def:g11:lithosphere-asthenosphere:lithosphere}{litosfer} muda
kehilangan banyak panas.

\textbf{4.} $490/\sqrt{4} = \qty{245}{mW/m^2}$.

\textbf{5.} $2.72 \times 0.026 \approx \qty{0.0707}{W/m^2} =
\qty{70.7}{mW/m^2}$: tepat di bawah rata-rata, praktis sama dengannya.

\textbf{6.} $0.0709 \times \num{2.073e14} \approx \qty{14.7}{TW}$.

\textbf{7.} $0.1054 \times \num{3.028e14} \approx \qty{31.9}{TW}$.

\textbf{8.} $14.7 + 31.9 \approx \qty{46.6}{TW}$, yaitu \qty{47}{TW} dari
perkiraan yang telah diterbitkan (rata-rata global dalam bab ini memberikan
\qty{46.7}{TW}, nilai yang dipakai di bawah).

\textbf{9.} $\num{4.66e13}/\num{5.101e14} \approx \qty{0.0914}{W/m^2}$,
sekitar \qty{91}{mW/m^2}.

\textbf{10.} $0.0243 \times \num{2.073e14} \approx \qty{5.0}{TW}$, sepertiga
dari \qty{14.7}{TW} yang hilang melalui benua: radioaktivitas \omterm{def:g12:continental-crust:continental-crust}{kerak benua}
merupakan bagian besar dari \omterm{def:g12:internal-heat:heat-flow}{aliran panas} benua.

\textbf{11.} Untuk setiap kilogram uranium, Bumi mengandung \qty{3.9}{kg}
torium, yang melepaskan $3.9 \times 26.3 \approx \qty{102.6}{\micro W}$,
dibandingkan dengan \qty{94.4}{\micro W} untuk uranium-238. Jadi uranium-238
menghasilkan $94.4/(94.4 + 102.6) \approx 0.479$ dari \qty{20.0}{TW}: sekitar
\qty{9.6}{TW}. Uranium-235: $9.59 \times 4.10/94.4 \approx \qty{0.42}{TW}$.

\textbf{12.} $20.0 + 0.42 + 4 \approx \qty{24.4}{TW}$.

\textbf{13.} Granit dan batuan lain \omterm{def:g12:continental-crust:continental-crust}{kerak benua}, sekitar dua puluh kali lebih
kaya daripada batuan ultramafik mantel.

\textbf{14.} Lebih besar: setiap isotop lebih melimpah pada waktu itu, dua kali
lebih banyak untuk setiap waktu paruh ke belakang; 3 miliar tahun yang lalu
jumlahnya lebih dari dua kali jumlah saat ini.

\textbf{15.} $46.7 - 24.4 = \qty{22.3}{TW}$.

\textbf{16.} Dari panas yang tersimpan sejak Bumi terbentuk: \omterm{def:g12:internal-heat:secular-cooling}{pendinginan sekuler}
mantel dan inti.

\textbf{17.} $22.3/46.7 \approx \qty{48}{\percent}$.

\textbf{18.} Dari $11/46.7 \approx \qty{24}{\percent}$ sampai $38/46.7 \approx
\qty{81}{\percent}$: geoneutrino sangat mempersempit rentang ini.

\textbf{19.} \omterm{def:g12:internal-heat:convection}{Konveksi mantel}, yang mengangkut sebagian besar panas; hanya
beberapa puluh kilometer terakhir yang dilintasi secara konduksi.

\textbf{20.} $24.4/46.7 \approx 0.52$: radioaktivitas memasok
\textbf{sekitar separuh} kehilangan panas Bumi saat ini, dan pendinginan
planet yang lambat separuh sisanya.
\end{solution}
